\chapter{Functions}

\chapterintro{A function assigns exactly one output to every input, which is the same promise a pure function makes in code, and it is the promise that makes caching, testing and composition safe. This chapter makes the idea precise and gives you composition and inverses, which the rest of the curriculum uses constantly.}

\section{Core Ideas}

\paragraph{Definition} A function is a rule that gives each input exactly one output. The domain is the set of allowed inputs; the range is the set of outputs actually produced. A relation that sends one input to two different outputs is not a function, which is what the vertical line test checks on a graph.

\paragraph{Notation} $f(x) = 3x - 2$ names the function $f$ and its rule. $f(4)$ means ``the output when the input is 4,'' which is $10$. $f(a + 1)$ means substitute $a + 1$ for every $x$.

\paragraph{Four representations} A table of input-output pairs, a graph, a formula, and a program. All four describe the same object; being able to move between them is the skill.

\paragraph{Domain restrictions} Division by zero and square roots of negatives are undefined in the reals, so $f(x) = 1/(x - 3)$ has domain all $x$ except $3$, and $g(x) = \sqrt{x - 4}$ has domain $x \ge 4$. In code these are the inputs that raise exceptions.

\paragraph{Composition} $(f \circ g)(x) = f(g(x))$: apply $g$ first, then $f$. Order matters, and $f(g(x))$ and $g(f(x))$ are usually different functions. Composition is what a pipeline of transformations is.

\paragraph{Inverses} The inverse $f^{-1}$ undoes $f$, so $f(f^{-1}(x)) = x$ and $f^{-1}(f(x)) = x$. To find it, swap $x$ and $y$ in the formula and solve for $y$. A function has an inverse only when no two inputs share an output.

\paragraph{Piecewise functions} Different rules on different parts of the domain. In code this is an if-else chain, and the conditions must not overlap.

\section{Worked Examples}

\begin{example}
Let $f(x) = 3x - 2$ and $g(x) = x^2$. Find $f(g(2))$ and $g(f(2))$.
\solution $f(g(2)) = f(4) = 10$, while $g(f(2)) = g(4) = 16$.
\end{example}

\begin{example}
Find the domain of $f(x) = 1/(x - 3)$.
\solution The denominator is zero when $x = 3$, so the domain is every real number except $3$.
\end{example}

\begin{example}
Shipping costs \$5 for the first pound and \$2 for each additional pound. Write the cost function and find the cost of a 3.5-pound package.
\solution $c(w) = 5$ if $w \le 1$, and $c(w) = 5 + 2(w - 1)$ if $w > 1$. For a 3.5-pound package, $c(3.5) = 5 + 2(2.5) = 10$.
\end{example}

\section{Exercises}

\begin{multicols}{2}
\begin{exercises}
\item For $f(x) = 2x^2 - x + 1$, find $f(0)$, $f(-2)$ and $f(a + 1)$, simplified.
\item Which of these are functions: (a) $\{(1, 2)$, $(2, 3)$, $(1, 4)\}$; (b) $\{(1, 2)$, $(2, 2)$, $(3, 2)\}$; (c) $y^2 = x$?
\item Find the domain of $g(x) = \sqrt{x - 4}$, and of $h(x) = \dfrac{x + 1}{x^2 - 9}$.
\item For $f(x) = x + 5$ and $g(x) = 2x$, find $f(g(x))$ and $g(f(x))$.
\item For $f(x) = x^2 - 1$ and $g(x) = 3x$, find $(f \circ g)(2)$ and $(g \circ f)(2)$.
\item Find the inverse of $f(x) = 4x - 3$ and verify that $f(f^{-1}(x)) = x$.
\item Let $f(x) = x^2$ for $x < 0$ and $f(x) = 2x + 1$ for $x \ge 0$. Find $f(-3)$, $f(0)$ and $f(4)$.
\item Find the range of $f(x) = x^2 + 2$ and of $g(x) = -|x| + 3$.
\item Temperature in Fahrenheit is $F(C) = 1.8C + 32$, and a comfort score is $s(F) = -|F - 72| + 100$. Write $s$ as a function of $C$ and evaluate it at 20 degrees Celsius.
\item \textbf{(proof)} Explain why $f(x) = x^2$ has no inverse on all real numbers, and why it does have one when the domain is restricted to $x \ge 0$.
\end{exercises}
\end{multicols}

\section{Programming Exercises}

\begin{codeexercises}
\item Write \lstinline!compose(f, g)! that returns a new function computing $f(g(x))$. Test Exercises~10.4 and 10.5 with lambdas, and show that \lstinline!compose(f, g)! and \lstinline!compose(g, f)! differ.
\item Write \lstinline!with_domain(f, allowed)! that returns a function raising \lstinline!ValueError! when \lstinline!allowed(x)! is false and otherwise returns $f(x)$. Wrap the two functions from Exercise~10.3 and test an input inside and outside each domain.
\item Represent a piecewise function as a list of \lstinline!(condition, rule)! pairs and write \lstinline!piecewise(rules)! that returns a function applying the first matching rule. Encode Exercise~10.7 and confirm your three values. Add a check that raises if no rule matches.
\item Write \lstinline!is_inverse(f, f_inv, samples)! that returns true when $f(f_{\text{inv}}(x))$ and $f_{\text{inv}}(f(x))$ both equal $x$ within tolerance for every sample. Test Exercise~10.6, and show that it fails for $f(x) = x^2$ with $f_{\text{inv}}(x) = \sqrt{x}$ when the samples include negatives.
\item Write \lstinline!memoize(f)! that caches results in a dictionary. Explain in a comment why this is only correct because $f$ is a function in the mathematical sense, and give an example of a Python callable it would break.
\end{codeexercises}
